\chapter{Text creation — Sampling Strategy and Decoding}
\label{ch:generation}

In this chapter, we implement the process of generating text with the learned GPT model. It covers a variety of strategies, from simple greedy decoding to Top-K, Top-P, and Temperature sampling.

\section{Principle of autoregressive generation}

The creation of the language model is done in the\keyterm{Autoregressive (autoregressive)}method. That is, it generates one token at a time, adds that token back to the input, and repeats the process to generate the next token.

\begin{enumerate}
\item Start with prompt$[w\_1, \dots, w\_k]$.
\item Put it in the model to get the following token distribution$P(w\_{k+1} | w\_1, \dots, w\_k)$.
\item Sample$w\_{k+1}$from the distribution (or argmax).
\item Append to sequence:$[w\_1, \dots, w\_k, w\_{k+1}]$.
\item 2$\sim$4 iterations until the ending condition (<eos> or maximum length).
\end{enumerate}

\begin{figure}[htbp]
\centering
\begin{tikzpicture}[
  every node/.style={font=\small\sffamily},
  tok/.style={draw=inkblue, fill=inkblue!8, minimum width=0.7cm, minimum height=0.6cm},
  newtok/.style={draw=inkred, fill=inkred!10, minimum width=0.7cm, minimum height=0.6cm},
  model/.style={draw=inkgray, fill=inkgray!10, rounded corners=2pt, minimum width=2cm, minimum height=0.7cm}
]
  % step 1
  \node[tok] (a1) at (0, 0) {The};
  \node[tok, right=0.05cm of a1] (a2) {cat};
  \node[model, right=0.5cm of a2] (m1) {model};
  \node[newtok, right=0.5cm of m1] (n1) {sat};
  \node[inkgray, anchor=south, font=\scriptsize] at (1.5, 0.5) {step 1};

  % step 2
  \node[tok] (b1) at (0, -1.5) {The};
  \node[tok, right=0.05cm of b1] (b2) {cat};
  \node[tok, right=0.05cm of b2] (b3) {sat};
  \node[model, right=0.5cm of b3] (m2) {model};
  \node[newtok, right=0.5cm of m2] (n2) {on};
  \node[inkgray, anchor=south, font=\scriptsize] at (1.5, -1) {step 2};

  % step 3
  \node[tok] (c1) at (0, -3) {The};
  \node[tok, right=0.05cm of c1] (c2) {cat};
  \node[tok, right=0.05cm of c2] (c3) {sat};
  \node[tok, right=0.05cm of c3] (c4) {on};
  \node[model, right=0.5cm of c4] (m3) {model};
  \node[newtok, right=0.5cm of m3] (n3) {the};
  \node[inkgray, anchor=south, font=\scriptsize] at (1.5, -2.5) {step 3};

  \draw[-{Stealth[length=4pt]}, inkblue] (a2) -- (m1);
  \draw[-{Stealth[length=4pt]}, inkblue] (m1) -- (n1);
  \draw[-{Stealth[length=4pt]}, inkblue] (b3) -- (m2);
  \draw[-{Stealth[length=4pt]}, inkblue] (m2) -- (n2);
  \draw[-{Stealth[length=4pt]}, inkblue] (c4) -- (m3);
  \draw[-{Stealth[length=4pt]}, inkblue] (m3) -- (n3);
\end{tikzpicture}
\caption{Autoregressive generation process. At each step, insert previous tokens into the model and predict the next token..}
\label{fig:autoregressive}
\end{figure}

\section{Greedy sampling}

The simplest strategy is to select the token with the highest probability at each step.

\begin{lstlisting}[language=Python]
class GreedySampler:
    def __call__(self, logits: torch.Tensor) -> torch.Tensor:
        # logits: [batch, vocab_size] -> token_ids: [batch]
        return torch.argmax(logits, dim=-1)
\end{lstlisting}

Greedy is deterministic and fast, but tends to select the same tokens every time, resulting in monotonous text. ``The cat sat on the mat. The cat sat on the mat. It's easy to fall into repetition like ''The cat...''.

\begin{warnbox}[Greedy’s Trap: Repeat Loop]
Greedy decoding produces boring text by selecting only ``safe'' tokens. A more serious problem is that it falls into a repetitive loop. Same phenomenon as ``I went to the store to buy some store to buy some store to buy...''.

Reason: If a loop is formed where the model predicts ``store'' followed by ``to buy'' with a high probability, and then ``to buy'' followed by ``some store'' with a high probability, greedy cannot get out of it.
\end{warnbox}

\section{Temperature sampling}

Introduces a parameter that controls the ``sharpness'' of the distribution. Divide the logit by the temperature$T$and then softmax.

\begin{formula}[Temperature Sampling]
\[
P\_T(w) = \frac{\exp(\text{logit}(w) / T)}{\sum\_{w'} \exp(\text{logit}(w') / T)}
\]
\begin{itemize}
\item $T < 1$: Distribution becomes sharper (closer to greedy)
\item $T > 1$: Distribution becomes flatter (more random)
\item $T = 1$: Original distribution
\end{itemize}
\end{formula}

\begin{lstlisting}[language=Python]
class TemperatureSampler:
    def __init__(self, temperature: float = 1.0):
        assert temperature > 0
        self.temperature = temperature

    def __call__(self, logits):
        probs = F.softmax(logits / self.temperature, dim=-1)
        return torch.multinomial(probs, num_samples=1).squeeze(-1)
\end{lstlisting}

\begin{figure}[htbp]
\centering
\begin{tikzpicture}[
  scale=0.8,
  every node/.style={font=\small\sffamily}
]
  % T=0.5 (sharp)
  \begin{scope}[shift={(0, 0)}]
    \node[inkblue, font=\bfseries\sffamily] at (3, 3.5) {T=0.5 (sharp)};
    \foreach \i/\h in {0/0.3, 1/0.5, 2/3.0, 3/0.4, 4/0.2, 5/0.1, 6/0.05} {
      \draw[inkblue, fill=inkblue!50] (\i*0.5, 0) rectangle (\i*0.5+0.4, \h);
    }
    \node[inkgray, font=\scriptsize, anchor=north] at (1.75, -0.1) {token};
  \end{scope}

  % T=1.0 (original)
  \begin{scope}[shift={(5, 0)}]
    \node[inkblue, font=\bfseries\sffamily] at (3, 3.5) {T=1.0 (original)};
    \foreach \i/\h in {0/0.5, 1/0.8, 2/1.8, 3/1.2, 4/0.7, 5/0.4, 6/0.2} {
      \draw[inkblue, fill=inkblue!50] (\i*0.5, 0) rectangle (\i*0.5+0.4, \h);
    }
  \end{scope}

  % T=2.0 (flat)
  \begin{scope}[shift={(10, 0)}]
    \node[inkblue, font=\bfseries\sffamily] at (3, 3.5) {T=2.0 (flat)};
    \foreach \i/\h in {0/0.8, 1/0.9, 2/1.2, 3/1.0, 4/0.85, 5/0.7, 6/0.6} {
      \draw[inkblue, fill=inkblue!50] (\i*0.5, 0) rectangle (\i*0.5+0.4, \h);
    }
  \end{scope}
\end{tikzpicture}
\caption{TemperatureDistribution changes according to. TIf is low, the distribution becomes sharp. greedygetting closer to, TThe higher it is, the flatter it becomes, increasing the randomness..}
\label{fig:temperature}
\end{figure}

\section{Top-K sampling}

To prevent tokens with low probability from being selected, only the top$K$tokens are considered.

\begin{lstlisting}[language=Python]
class TopKSampler:
    def __init__(self, k: int = 50, temperature: float = 1.0):
        self.k = k
        self.temperature = temperature

    def __call__(self, logits):
        scaled = logits / self.temperature
        top_values, top_indices = torch.topk(scaled, self.k, dim=-1)
        probs = F.softmax(top_values, dim=-1)
        sampled = torch.multinomial(probs, num_samples=1).squeeze(-1)
        return top_indices.gather(-1, sampled.unsqueeze(-1)).squeeze(-1)
\end{lstlisting}

\subsection{Top-Kintuition}

Cut off the ``tails'' (tokens with very low probability) from the logit distribution. For example, when the vocabulary is 50,000, considering only the top 50 means sampling from tokens with a probability mass of 99.99\%.

If Top-K = 1, it is the same as greedy. If Top-K =$V$(vocabulary size), it is the same as temperature sampling.

\section{Top-P (Nucleus) Sampling}

Top-K fixes$K$, but Top-P dynamically adjusts the number of tokens considered depending on the distribution type. Only tokens until the cumulative probability reaches$P$are considered.

\begin{lstlisting}[language=Python]
class TopPSampler:
    def __init__(self, p: float = 0.9, temperature: float = 1.0):
        self.p = p
        self.temperature = temperature

    def __call__(self, logits):
        scaled = logits / self.temperature
        probs = F.softmax(scaled, dim=-1)
        sorted_probs, sorted_indices = torch.sort(probs, descending=True, dim=-1)
        cumulative = torch.cumsum(sorted_probs, dim=-1)
        # PMasking from positions beyond
        sorted_mask = cumulative - sorted_probs > self.p
        sorted_probs = sorted_probs.masked_fill(sorted_mask, 0.0)
        sorted_probs = sorted_probs / sorted_probs.sum(dim=-1, keepdim=True).clamp(min=1e-10)
        sampled = torch.multinomial(sorted_probs, num_samples=1).squeeze(-1)
        return sorted_indices.gather(-1, sampled.unsqueeze(-1)).squeeze(-1)
\end{lstlisting}

\begin{figure}[htbp]
\centering
\begin{tikzpicture}[
  scale=0.7,
  every node/.style={font=\small\sffamily}
]
  % Top-K visualization (left)
  \begin{scope}[shift={(0, 0)}]
    \node[inkblue, font=\bfseries] at (3, 4.5) {Top-K (K=5)};
    \foreach \i/\h/\sel in {0/2.0/1, 1/3.5/1, 2/1.5/1, 3/0.8/1, 4/2.8/1,
                             5/0.3/0, 6/0.2/0, 7/0.15/0, 8/0.1/0, 9/0.05/0} {
      \pgfmathsetmacro{\col}{ifthenelse(\sel == 1, "inkblue!50", "inkgray!20")}
      \draw[fill=\col, draw=inkblue] (\i*0.5, 0) rectangle (\i*0.5+0.4, \h);
    }
    \draw[inkred, dashed, line width=0.8pt] (2.5, 0) -- (2.5, 4);
    \node[inkred, font=\scriptsize, anchor=west] at (2.6, 3.8) {K=5};
  \end{scope}

  % Top-P visualization (right)
  \begin{scope}[shift={(8, 0)}]
    \node[inkblue, font=\bfseries] at (3, 4.5) {Top-P (P=0.9)};
    \foreach \i/\h/\sel in {0/2.0/1, 1/3.5/1, 2/1.5/1, 3/0.8/1, 4/2.8/1,
                             5/0.3/1, 6/0.2/0, 7/0.15/0, 8/0.1/0, 9/0.05/0} {
      \pgfmathsetmacro{\col}{ifthenelse(\sel == 1, "inkblue!50", "inkgray!20")}
      \draw[fill=\col, draw=inkblue] (\i*0.5, 0) rectangle (\i*0.5+0.4, \h);
    }
    \node[inkred, font=\scriptsize, anchor=west] at (3.0, 3.8) {cumsum};
    \node[inkred, font=\scriptsize, anchor=west] at (3.0, 3.4) {$\approx$ 0.9};
  \end{scope}

  % axis
  \draw[inkblue, ->] (-0.5, 0) -- (5.5, 0) node[right, font=\scriptsize] {token};
  \draw[inkblue, ->] (-0.5, 0) -- (-0.5, 4) node[above, font=\scriptsize] {prob};
  \draw[inkblue, ->] (5.5, 0) -- (10.5, 0) node[right, font=\scriptsize] {token};
\end{tikzpicture}
\caption{Top-K vs Top-P. Top-Kis a fixed number, Top-Pis the cumulative probability PSelect until.}
\label{fig:topk-topp}
\end{figure}

\subsection{Top-K vs Top-P}

\begin{center}
\begin{tabularx}{\textwidth}{lXX}
\toprule
 & \textbf{Top-K} & \textbf{Top-P (Nucleus)} \\
\midrule
Number of tokens considered & fixed ($K$) & dynamic (depending on distribution) \\
When distribution is sharp & consider too many tokens & only consider few tokens (automatic) \\
When distribution is flat & consider too few & consider many tokens (automatic) \\
Adaptability & Low & High \\
\bottomrule
\end{tabularx}
\end{center}

Top-P usually produces more natural results. This is because the model considers only a few tokens when it is confident, and many tokens when it is uncertain.

\section{Sampling Strategy Comparison}

\begin{center}
\begin{tabularx}{\textwidth}{lXXX}
\toprule
\textbf{strategy}&\textbf{manifold}&\textbf{consistency}&\textbf{example}\\
\midrule
Greedy & Low & High & Code Generation, Math \\
Temperature (T=0.7) & Medium & Medium & General use \\
Top-K (K=50) & Medium & High & Chatbot \\
Top-P (P=0.9) & Medium & High & Chatbot, Creation \\
Top-P (P=0.95) & High & Medium & Creative Writing \\
\bottomrule
\end{tabularx}
\end{center}

\section{Implementing the creation function}

Create a\code{generate}function that combines all samplers.

\begin{lstlisting}[language=Python]
@torch.no_grad()
def generate(model, tokenizer, prompt, max_new_tokens=50,
             sampler=None, device="cpu", stop_on_eos=True):
    if sampler is None:
        sampler = GreedySampler()
    model.eval()
    model.to(device)
    ids = tokenizer.encode(prompt, add_bos=True, add_eos=False)
    input_ids = torch.tensor([ids], dtype=torch.long, device=device)
    eos_id = tokenizer.special.eos_id
    context_length = model.config.context_length

    for _ in range(max_new_tokens):
        # Context length limit
        x = input_ids[:, -context_length:] if input_ids.shape[1] > context_length else input_ids
        logits = model(x)
        next_logits = logits[:, -1, :]  # last location
        next_id = sampler(next_logits)
        input_ids = torch.cat([input_ids, next_id.unsqueeze(0)], dim=1)
        if stop_on_eos and next_id.item() == eos_id:
            break
    return tokenizer.decode(input_ids[0].tolist())
\end{lstlisting}

\section{Context length limit}

The model can only process up to\code{context\_length}specified during training. During construction, if the sequence exceeds this length, the oldest token is truncated (sliding window). It's an approximation, but it works. Volume 2 covers RoPE's extrapolation technique to overcome this limitation.

\begin{warnbox}[Problem with Sliding Window]
If you cut off the context, the model will forget the ``front part''. ``Once upon a time, there was a princess named Aurora. She lived in a castle. One day, [Long explanation] ... If you cut out the first part of 'Aurora decided to'', the context that ``Aurora'' is the princess's name is lost.

Solution: Use a longer context model (RoPE extrapolation), or compress it by summarizing it. Covered in Volume 2.
\end{warnbox}

\section{sampler test}

\begin{lstlisting}[language=Python]
def test_greedy_picks_max():
    sampler = GreedySampler()
    logits = torch.tensor([[1.0, 3.0, 2.0, 0.5]])
    out = sampler(logits)
    assert out.item() == 1  # largest value

def test_top_k_limits_choices():
    torch.manual_seed(0)
    logits = torch.tensor([[1.0, 5.0, 4.0, 3.0, 0.1]])
    sampler = TopKSampler(k=2, temperature=1.0)
    for _ in range(20):
        idx = sampler(logits).item()
        assert idx in (1, 2)  # always top-one of 2

def test_top_p_limits_choices():
    torch.manual_seed(0)
    logits = torch.tensor([[0.0, 10.0, 9.0, 0.0, 0.0]])
    sampler = TopPSampler(p=0.5, temperature=1.0)
    for _ in range(20):
        idx = sampler(logits).item()
        assert idx == 1  # Only the highest probability tokens
\end{lstlisting}

\section{practice: Comparison of various sampling}

\code{scripts/tiny\_train.py}Running the demo allows you to compare the output of five different sampling strategies.

\begin{lstlisting}[style=bashstyle]
cd make-llm-basic
python scripts/tiny_train.py --epochs 5 --vocab 500
\end{lstlisting}

Example output:
\begin{lstlisting}[basicstyle=\ttfamily\footnotesize, backgroundcolor=\color{codebg}, frame=single]
prompt: 'the'
  [Greedy            ] the cat sat on the mat
  [Temperature=0.5   ] the lazy dog sleeps
  [Temperature=1.0   ] the brown fox jumps
  [Top-K=10          ] the quick fox runs
  [Top-P=0.9         ] the warm sun shines
\end{lstlisting}

\section{Advanced decoding techniques (introduction)}

Although not implemented in this book, we introduce advanced techniques used in real systems:

\subsection{Beam Search}
At each step,$K$candidate sequences are maintained, and the sequence with the highest score is selected at the end. Suitable for translation, summarization, etc., but not suitable for chatbots (lack of variety).

\subsection{Speculative Decoding}
The smaller ``draft'' model generates$K$tokens first, then the larger model verifies it in parallel. If correct, adopt$K$at once, if incorrect, start from the beginning. Throughput 2$\sim$can be improved by 3 times.

\subsection{Contrastive Search}
Modern techniques for maintaining consistency while avoiding repetition. Select various tokens by penalizing them for similarity to previous token representations.

\section{summation}

\begin{itemize}
\item Autoregressive generation generates one token at a time and extends the sequence.
\item Greedy: The token with the highest probability. Dangers of monotony and repetitive loops.
\item Temperature: Controls the sharpness of the distribution. The lower T, the more deterministic.
\item Top-K: Only the top K items are considered. Simple but non-adaptive.
\item Top-P: Only tokens up to cumulative probability P are considered. Adaptive to distribution.
\item Truncates the oldest token if the context length is exceeded.
\item Advanced techniques such as speculative decoding and beam search also exist.
\end{itemize}

\section{practice problems}

\begin{enumerate}
\item Analyze why greedy decoding falls into an iterative loop. Under what kind of distribution do loops occur most often?
\item How is it different from greedy when Temperature approaches 0? What happens if it is 0?
\item What is the effect of applying Top-K and Top-P at the same time? (i.e. filtering Top-P among Top-K candidates)
\item What sampling is best for chatbots? Explain why.
\item Why is speculative decoding fast?
\end{enumerate}

This completes the core implementation of Volume 1. The next chapter summarizes the entire book and introduces advanced topics that will be covered in Volume 2.
